\section{Method --- 모듈별 생각}
\label{sec:method}

\subsection{두 모델 제약}
\begin{itemize}
  \item Jev: 텍스트만. 수치·시간 비교 약함. 입력 길면 정확도 하락. \textbf{Score 기댓값으로 크기 보간 금지}, \textbf{질문끼리 확률이 서로 안 맞음}(공식 문서)
  \item Astra: 이미지 OK, 비디오 없음. \textbf{logprob 못 씀}(effort none 없음). effort 지연은 날마다 바뀜
  \item \chg{(09-24 §44) Jev 사용 불가 $\rightarrow$ Jev-L(로컬 VLM, 이미지 받음). 위 Jev 줄의 ``Score 보간 금지·질문끼리 확률 안 맞음''은 Jev 문서 근거였지만 설계 규칙으로 유지(기하는 코드, 같은 질문 보기끼리만 확률 비교). 모델 고정 = 저장소 리비전 + vLLM 판본 + BF16(\texttt{jev-1.13.0} 대신), \texttt{JevCall} $\rightarrow$ \texttt{DecCall}}
\end{itemize}

\subsection{M1 이미지 $\rightarrow$ Jev}
\begin{itemize}
  \item \intent{Jev는 이미지를 못 본다. LLM/VLM/VLA 논문뿐 아니라 텍스트와 이미지를 제대로 변환하는 다른 분야 연구에서 가장 효율적인 방법을 찾는다.}
  \item \idea{형식보다 내용. 작업 관련 물체만 + 코드가 계산한 관계·술어}
  \item \chg{``텍스트가 이미지보다 낫다''의 근거(BALROG)는 모델마다 다름. 근거에서 내림}
  \item \chg{(09-24) 변환 방법은 Claude가 정함: 기본 = 물체 ID 술어 표 + 변화 절(후보 A). 나머지 두 후보는 오프라인 비교(E3) 조건. ``변환 방법은 사용자가 정한다는 사용자가 한 말이 아니라서 지움} \decide{인식 앞단}
  \item \idea{(09-24 D13) 이미지를 못 보는 모델도 기호 표만 받아 게임 벤치 상위권(lmgame-Bench, 저자도 예비 결과라 함). 가장 강한 VLM(GPT-4.1·GPT-5.2)도 ``옆에 있음·열림'' 판정은 50--77\%(ViPlan~\cite{viplan2025}) $\rightarrow$ 그런 술어는 코드로 계산}
  \item \chg{(09-24 §44) Jev-L은 이미지를 받지만 작은 VLM의 공간 판단은 못 믿음 $\rightarrow$ M1 텍스트가 주 입력, 머리캠 1장은 보조(텍스트만 조건과 비교)}
  \item \idea{(E3-ST, fx 367 재계산) 분할기를 바꿔도 정지 흔들림은 거의 같음(같은 에피소드 과제 1: 1.45 대 1.42 mm, 과제 2: 2.41 대 2.47 mm). 성패는 검출 프롬프트: SAM 3.1에 과제 문장 ``box''는 흰 바구니 0/20, ``basket''은 19/20. 정지 p95 9--16 mm라 과제 2는 near 띠(1 cm)를 넘음}
  \item \idea{(Claude 보고, user-log 49) 물체 이름은 Astra가 계획할 때 장면을 보고 정해 주는 구조가 맞아 보입니다.}
  \item \intent{나도 이렇게 생각함}
  \item \chg{(09-24 §46) SAM 3.1 프롬프트 = Astra가 계약 \texttt{objects}에 적은 \texttt{detect\_phrase}. 검출 0건·추적 끊김이면 \texttt{id\_uncertain}(술어 unknown, 소프트 채널)로 올리고 다음 하트비트를 앞당겨 Astra가 이름을 고침(정지 없음). T0에서는 검출 확인 뒤 실행 시작. 계약 검사기에 ``detect\_phrase 비어 있지 않음'' 추가}
  \item \chg{(09-24 §47) ZED Mini ``102$^\circ$''는 센서 최대 화각. 우리가 쓰는 VGA 672$\times$376은 수평 85$^\circ$, fx 367(실측 364.0) $\rightarrow$ 깊이·흔들림 다시 계산(깊이 1.349배, 흔들림 1.15--1.30배)}
\end{itemize}

\subsection{M2 Astra $\rightarrow$ Jev}
\begin{itemize}
  \item \intent{1번과 거의 같은 문제다. Astra는 이미지를 보고 명령하지만 Jev는 아니니, 효율적으로 전달하는 방법을 더 생각한다.}
  \item \idea{(``1번'' = M1) 물체 번호로 말하게 함(텍스트판 Set-of-Mark, 기간 밖, 기초 문헌). 제약은 Astra가 코드로 쓰고 Jev에는 참/거짓만}
  \item \idea{작은 모델에겐 ``어떻게''보다 ``무엇을''(목표형)이 나음. 계획이 본 시각·금지 조건·단계 ID 칸 추가}
  \item \chg{(09-24 D13) ``작은 실행기엔 목표만'' 근거였던 COPE는 실제로 큰 모델의 짧은 \emph{절차} 계획이 도왔음 $\rightarrow$ 근거에서 뺌. 목표 중심은 설계 선택으로 두고 ``목표 + 절차 1--2문장'' 조건으로 확인}
  \item \chg{(09-24 D18) Astra 계획 확인: 3번 다수결은 삭제 $\rightarrow$ 한 번 부르고 코드로 검사(전제·효과 포함), 틀리면 위반 목록 붙여 최대 2번 고치게 함. 실행 중엔 이전 계약 유지하고 뒤에서 고침(선행: 계획 검증 반복 수리, JEV-Star)}
  \item \chg{(09-24 D18) 재계획 형식: patch 고정 $\rightarrow$ 기본 patch, 전제가 크게 깨지거나 목표가 바뀌거나 patch가 검사 불통과면 전체 재작성. 실행한 단계는 불변(선행: Plan-and-Act, SWE-Edit 적응형)}
\end{itemize}

\subsection{M3 행동}
\begin{itemize}
  \item \intent{촘촘하게 가도 된다. Jev 방식이면 충분히 가능하고 할 만하다고 본다.}
  \item \chg{``확률 평균으로 크기''는 Jev 공식 문서가 금지 $\rightarrow$ 이름 붙은 구간을 고름}
  \item \chg{``보기 255개 한 번이면 충분'' 폐기 $\rightarrow$ 질문당 서로 뚜렷한 보기 $\le$ 약 17개. D255는 비교 조건만}
  \item \chg{(09-24) 보기 이름: 2지선다는 중립 ID(A/B) + 설명, 3개 이상은 내용 이름(+x, 컵 등) + 설명. 3개 이상 쪽은 E0.5로 확정}
  \item \chg{(09-24 D18) E1 뒤 켤 후보: 질문별 conformal 집합. 하나면 확정 후보, 둘 이상이거나 모름 포함이면 보류(감속), 반복되면 Astra로(선행: KnowNo, CoFineLLM)}
  \item \idea{반대 증거 있음: 학습 없는 모델은 잘게 여러 번 움직이기보다 목표 지점을 한 번에 가리킬 때 훨씬 잘함(드론: 연속 수치 글 7\%, 후보 점 선택 40\%, 목표 점 직접 지정 87--100\%~\cite{spf2025}, GUI 반복 이동 붕괴~\cite{guicursor2025})}
  \item \dirn{``목표 지정형''(Jev가 목표 물체·부분을 고르고 크기는 코드)을 원안 옆에 나란히 재 봄. \decide{기본안}}
  \item \idea{RoboDawn(VLM이 cm 숫자로 로봇 조종)의 하네스 구성 요소 절제(Gemini-3.8-Flash 0-shot)에서 가장 크게 움직인 건 격자 같은 공간 근거(빼면 47.0$\rightarrow$32.4)이고 결정 횟수도 움직임(샷 수·모델 선택은 이보다 큼) $\rightarrow$ 촘촘함은 결정 빈도·공간 근거로 살림. 저자도 ``이산 명령은 미세 보정에 거칠다''고 함 $\rightarrow$ 접촉 근처 미세 조정}
  \item \dirn{첫 실험 = E0 지연 $\rightarrow$ E0.5 표 재생 $\rightarrow$ E1 보정 $\rightarrow$ E2a 마차(같은 텍스트 상태에서 룰 대 Jev 촘촘·목표 지정·혼합, M4 없이). 겹침 켬·끔은 E-M4에서. 방향 재검토는 95\% 상한이 +5\%p 미만일 때만}
\end{itemize}

\subsection{M4 RTC식 겹침 호출 (가장 중요)}
\begin{itemize}
  \item \intent{1초에 Jev를 한 번이 아니라 3번 정도 계단식으로 겹쳐 쓴다.}
  \item \intent{겹치는 부분의 움직임이 제대로 반영됐는지 스스로 감지하고 업데이트하는 시스템이 자동으로 있어야 한다. 이게 굉장히 중요하다.}
  \item \intent{Jev가 빠르다지만 로봇 동작에서는 빠르지 않으니, 연속 동작처럼 내뱉게 하는 법을 RTC에서 많이 참고한다.}
  \item \idea{RTC~\cite{black2025rtc}의 고정·중간·끝 구간을 이산판으로 가져옴(아래 방향 줄)}
  \item \idea{확인을 둘로 나눔: (a) 호출끼리 같은 말을 하나(앞부분 합의) (b) 실제로 그렇게 움직였나(예상 상태 대 측정 상태, 코드)}
  \item \idea{(a)만으론 안 됨. 같은 모델은 같은 오답을 반복}
  \item \idea{가장 가까운 선행 A3 저자가 스스로 ``정책 내부 신호만으론 자신 있게 틀린 예측을 못 잡는다, 외부 증거가 향후 과제''라고 씀 $\rightarrow$ 우리 (b)가 바로 그 자리}
  \item \chg{``안 멈추고 겹쳐 보내기''는 선행 있음(Slow Brain~\cite{slowbrain2026}: VLM Stream = 최신 답만, Probability Fusion = 최신 답 + 빠른 층 점수 감쇠 융합. Jev 데모: 요청 하나 + 직전 행동 유지). 합의로 앞 구간 확정도 A3가 먼저 함 $\rightarrow$ 새로움은 시간차 호출 합의 + (b) 결합으로 전제 무효화·keep/replace/repair}
  \item \dirn{E-M4 조건(E2a와 같은 단일 pick-and-place 과제): C0 멈춤 / C1 요청 하나 / C2 Slow Brain VLM Stream / C2' Slow Brain Probability Fusion(원문판) / C2'' 우리 변형(표 나이 감쇠 가중 최빈) / C5-A3(같은 시각 여러 샘플 합의만, (b) 없음) / C-FIX(합의 없이 앞 $k$ 스텝 고정 확정, $k$ 격자 탐색) / C3 (a)만 / C3' 베이즈 정지 / C3'' 보기 순서 돌리기 / C4 (b)만 / C5 (a)+(b) / C5' (b) 범주를 Jev에 안 줌 / C6 겹침 끔(측정용)}
  \item \dirn{핵심 비교: C5 대 C2·C2', C5 대 C5-A3, C5 대 C-FIX. 추가 E-M4-lat: 인공 지연 p95 0.3/0.8/2/5초 스윕} \chg{(09-24) 논문 한 편이라 지연 스윕과 RoboDojo 새 장면 실시간 트랙 M4 켬/끔(E-link)이 필수}
  \item \chg{VLASH~\cite{vlash2025}·A2C2~\cite{a2c2}는 학습형이라 발상만. A2C2는 ICLR 탈락}
  \item \dirn{RTC 3구간의 이산판: 고정(길이 = 실측 Jev p95) / 중간(바꾸려면 합의가 필요하고, 고정 경계에 가까울수록 문턱이 엄격 --- 가중 평균은 안 씀) / 끝(가장 새 표를 가확정)}
\end{itemize}

\subsection{M5 스무딩}
\begin{itemize}
  \item \intent{두 단계: 먼저 RTC와 유사한 방법, 다음으로 액션 스무딩이나 VLA에서 쓰는 부드럽게 하는 방법으로 아주 작은 움직임을 처리한다.}
  \item \idea{겹침 평균은 정밀 작업 성공률을 깎음 $\rightarrow$ 같은 보기끼리만 RTC식으로 섞고, 보기가 바뀌면 섞지 않고 관성화(접촉 근처 끔), 마지막에 저크 제한(Ruckig, 잠정). LiPo~\cite{lipo2025}는 비교}
  \item \chg{(09-24 D13) RTC 원문 가중은 끝에서 정확히 0이 되는 모양 --- 우리 지수식은 블렌드 끝에 약 5\% 계단이 남아 원문 모양으로 바꿈. 블렌딩 기본 유지는 그대로}
  \item \idea{WAM(심사 전, 조건당 5회, 신뢰도 낮음)의 ``사후 블렌딩은 이득 없음''은 한 과제만의 결과였음(다른 과제는 20$\rightarrow$40, 44$\rightarrow$80) $\rightarrow$ 근거에서 뺌. 가정이 바뀐 직후엔 블렌딩 끄는 조건을 비교}
\end{itemize}

\subsection{M6 스킬 $\leftrightarrow$ Jev}
\begin{itemize}
  \item \intent{기존 스킬도 Jev와 엄청 연관될수록 좋다.}
  \item \intent{1년 이내에 나온 GitHub 스타 상위 5개 중에서, LLM이 스킬을 불러오고 끝이 아니라 잘게 엮어 놓은 논문을 참고해 파이프라인을 대충 짜 본다.}
  \item \idea{스타 상위: PhyAgentOS, Zetta, Harness VLA, CaP-X, ROSClaw(저장소, 논문 없음. 또는 Show-Harness)}
  \item \idea{스타 1위는 거칠게 엮음. 잘게 엮는 쪽은 Zetta~\cite{zetta2026}(런타임 critic은 코드. LLM 승인자가 온라인인지는 원문이 엇갈림)와 Show-Harness}
  \item \dirn{(09-24) Harness VLA~\cite{harnessvla2026} 코드에서: 스킬 성공 여부 한 비트는 틀림(``성공=거짓이어도 영상상 들고 있으면 계속''이 실제 교훈) $\rightarrow$ 스킬 출력에 원자료 진단값(들어올린 높이·그리퍼 간격 등)을 필수로. 결정 지점 후보(자세 잡기·지금 호출·청크 예산·재시도 문장·잡힘 확인)는 원문 탐색 레버에서 가져옴}
  \item \idea{CaP-X~\cite{fu2026capx}의 생성 스킬 9개는 거의 좌표 계산 함수, 조건·자동 검증 없음 $\rightarrow$ 생성 스킬(b)엔 계약 + 실행 기반 검증을 붙임. 사람이 만든 높은 단계 API일수록 성공률이 오름 $\rightarrow$ ``무엇 위에서 생성하나''를 실험 축으로}
  \item \idea{Harness VLA 저자 스스로 ``계획기와 하위 사이 루프가 열려 있다''고 함 $\rightarrow$ 우리 M4/M7 자리}
  \item \idea{상위 6편은 상위 모델 결정 중 로봇이 기다리거나 상위 모델이 루프 밖. 다만 안 멈추고 겹쳐 부르는 것 자체는 Slow Brain~\cite{slowbrain2026}·Jev 데모에 있음 $\rightarrow$ 우리 자리는 시간차 호출 합의 + 실행 뒤 예상 대 측정을 결합한 확정 규칙(M4, 합의 확정 자체는 A3에 있음)과 스킬 안 결정 지점(M6)}
  \item \dirn{대략 파이프라인: Astra가 ``세션 계약''(단계 원장·스킬·결정 지점·critic 규칙) $\rightarrow$ Jev 3Hz 겹침 결정 $\rightarrow$ 스킬 100Hz $\rightarrow$ 코드 critic $\rightarrow$ 복구 사다리 $\rightarrow$ 검증된 것만 기억}
  \item \chg{(09-24) 결정 지점 보기도 같은 규칙: 두 갈래는 중립 ID + 설명, 3개 이상은 내용 이름 + 설명(E0.5로 확정)}
  \item \decide{5위 ROSClaw 대 Show-Harness}
  \item \intent{액션이 스포트하고 같은 비슷한 학습 모듈이 추가되어 있으면 좋을 것 같은 생각도 드는데 액션 익스펙트 스킬들을 통한 GT값이 많으면 해결이 될것 같은데}
  \item \intent{스킬들의 모든 실행이 거기들어가 있어야한단뜻}
  \item \intent{가능하면 엔드투엔드 혹은 다른모듈도 포함시키는 소규모엔드투엔드 학습 중 올바른거}
  \item \chg{(09-24 사용자 결정 B) 스킬 실행은 스킬 GT로 학습한 학습 실행기 하나가 맡음(스크립트 스킬 = GT 생성기 + 백업). Astra·Jev 결정 층은 학습 안 함. 계약·진단·FAIL은 코드. 주 표 실행기는 스크립트(잠정), 학습 실행기 표는 병기}
  \item \intent{무작위, 스킬로해도 좋긴한데 엑션 익스퍼트가 역할을 부가할방법}
  \item \intent{Jev는 앞으로, 뒤로 좌우 몇 센티를 판단하는 역할임 알지?}
  \item \chg{(09-24, 위 결정 B를 바꿈) 도메인 무작위화 넣음, 주 표 실행기는 스크립트 스킬. 학습 모듈은 스킬 명령 위에 작은 잔차 보정만 더함(Jev가 고른 방향·크기 구간 안에서만, 코드로 자름, 접근·접촉 근처에서만, 이미지 없음). 옆에서만 돌리는 진행 예측은 내부 절제로만. 학습 실행기는 비교 표로만}
\end{itemize}

\subsection{M7 진행 판별}
\begin{itemize}
  \item \intent{프레임별로 확인해서 Astra의 계획대로 되고 있는지 판별한다.}
  \item \intent{VLA에서 ``계획대로''를 무엇으로 정의했는지 많이 참고한다. VLA는 로봇을 학술적으로 조종하는 개념 공부를 많이 해 놓았다.}
  \item \idea{v2 그대로(센서·술어, 남은 시간). 시간 임계값은 SAFE~\cite{safe2025}의 conformal 보정에서 착안, 가끔 Astra로 GVL~\cite{gvl2024}(기간 밖, 기초 문헌)}
  \item \idea{실패 판정은 코드 critic 한 곳에서만(정체 규칙·마감 초과·겹침 불일치·Jev 진행 범주를 모아서)}
  \item \chg{$\pi^*_{0.6}$·$\pi_{0.7}$의 진행 정의는 ``계획대로 정의에 흡수하되, $\pi$ 계열이 실행 중 벗어남을 판정한다고는 안 씀(학습 신호) $\rightarrow$ 술어+남은 시간 유지}
  \item \chg{(09-24) ``적시 재현율''을 CheckVLA~\cite{checkvla2026} 방식이라 적었던 건 틀림 $\rightarrow$ 원문은 되돌릴 수 없는 시점 기준. 우리 고정 창(2초)판과 원문판을 같이 잼. 오경보는 에피소드 단위}
  \item \dirn{판정 기준은 고정 창 재현율($\tau = 2$ s)로 읽음. CheckVLA 원문판은 병기 보고만(Astra 지연이 날마다 크게 달라 사전 등록 기준으로 불안정)}
  \item \dirn{정체 규칙 원문판(Critic in the Loop~\cite{criticloop2026}: 진행 최고점이 약 9초(우리 계산) 갱신 안 되면) 비교 조건 추가, 우리 2초와 스윕}
\end{itemize}

\subsection{M8 Astra 호출}
\begin{itemize}
  \item \intent{적어도 실패 판정이 났을 때는 Astra를 부른다. 그 순간 한 장이 아니라 연속 프레임을 준다. 예: 이미지 하나에 사진 10개.}
  \item \intent{호출 방식은 비교 평가한다: Hz 단위 주기 / 정해 둔 필요한 순간마다 / ``이때쯤 결과가 나왔어야 하는데 안 나오네''라는 자체 판단.}
  \item \intent{Astra가 처음 생각할 때는 멈추는 게 맞지만, 그 이후는 웬만하면 멈추지 않는다. 실패로 생각해야 하는 경우만 어쩔 수 없다.}
  \item \idea{``이때쯤 나왔어야 하는데''의 구현 두 가지를 비교: 코드가 예상 시간으로 마감 계산 / 모델이 경과 시간 보고 스스로 판단 \decide{}}
  \item \chg{v3에서 기본값을 개별 이미지로 바꿨던 것 되돌림 $\rightarrow$ 사용자 예시(격자 한 장에 10프레임)가 기본 조건, 개별 이미지·변화 텍스트는 비교 조건. 광류 키프레임 선택과 덧그림의 근거~\cite{kite2026}}
  \item \intent{애플트는 하이로 둘다 비교를 해보면 좋을 것 같아}
  \item \intent{D14의 경우에는 아 안되겠다 애폴트 하이한 경우엔 73초 걸리는 거면은 로우로 가야겠다 그건 너무 느린 것 같아}
  \item \chg{(09-24 사용자 결정) effort 기본 low(high는 첫 토큰 약 73초라 너무 느림, low는 약 3초). 실험은 low와 high 둘 다 비교(medium·xhigh·max는 여력 있으면). 예전 ``잠정 기본 high''와 다른 작업 발언 ``추론 강도 낮추지 말 것''은 이 프로젝트에 적용 안 함}
  \item \idea{effort는 성공률보다 결정 수·시간을 바꿈(Show-Harness 스텝 37$\rightarrow$30, GPT-5.6-sol, 프로젝트 페이지). 지표에 결정 수 추가. 단 RPent Astra low 92.63\%라는 강한 값도 있음(논문 아닌 저장소 리더보드 값, LIBERO-PRO T+S 8칸, seed 1--10, 메모리 스냅샷 157/200 + 584/600, 에피소드당 약 412 s 멈춰서 생각) $\rightarrow$ \chg{(09-24) 기본 low로 해소, low·high 둘 다 비교}}
  \item \dirn{프레임 비교에 ``무엇이 바뀌었는지 텍스트'' 조건 추가(CaP-X에서 원시 이미지보다 나음)}
  \item \dirn{(09-24) 비교 기준을 원문 그대로: 호출 수 맞춘 주기(CheckVLA), 정체 때 멈추고 부르는 판(Critic in the Loop). ``실패 트리거를 다음 주기까지 미루기''는 절제로만 --- 사용자 원칙은 그대로}
  \item \idea{반대 증거: CheckVLA는 지연 1초에서 이미 ``다음 경계까지 기다림''(34.0)이 청크 안 수리(23.4--31.4)보다 나았음(Table S6에서 우리가 읽은 값). Astra는 훨씬 느림 $\rightarrow$ 그동안 아래 층이 버텨야 함(M9)}
  \item \intent{근데 처음 계획 세우고 실패할 때만 불러온다는 것 자체도 약간 좀 이상한데 어느 정도 주기로 하든지 뭔가가 있어야 될것 같은데 제미나이 3.0에서 어떤 식으로 했는지 한번 찾아봐봐 상위 계획기를 어떤 식으로 사용했는지 VLA 할때}
  \item \idea{(D25) Gemini Robotics-ER 문서~\cite{geminier_docs2026}: 최신 프레임 + 짧은 질문을 주기적으로 보내 명시적 결정을 강제하는 하트비트, 턴 완료를 기다리며 약 1 Hz 목표, 타이머만 믿으면 취소 반복 고리 경고. Gemini Robotics 1.5~\cite{geminirobotics15_2025}: 하위 과제 성공 감지로 다음 단계. Hi Robot 1 s 또는 사용자 개입 때(기간 밖, 기초 문헌) $\rightarrow$ 배포형 시스템은 주기 확인 + 단계 경계 + 사건을 섞음}
  \item \chg{(09-24 §45) H-cadence: T0(정지 허용) + T\_fail(항상) + 하트비트 T\_hb(직전 응답 뒤 $N$초, ack/patch/replace 닫힌 선택, 한 번에 1개, 잠정 $N$ = 5 s, 실효 약 8--9 s) + 단계 경계 T\_sub(기다리지 않음) + 사건(C\_assume·T3a·T\_stag·T\_j5·새 지시)은 하트비트를 앞당김. ack는 전제 판본을 안 올림(M4 표 보호). $N$은 E-M8c(K0--K4, 감시기 밖 실패 섭동 필수)로 정함}
  \item \idea{(09-24 실측) Astra low 첫 토큰 2.975 s(1회, 짧은 입력). 설계값 아님}
  \item (5 s 하트비트·단계 경계 확인 강등 질문에 답) \intent{바꾸기 (Recommended)}
  \item (오프라인 일 effort 질문에 답) \intent{논문에 로우 하이를 모두 가져가는 쪽이 맞겠네 성능비교두 그렇고}
  \item (Astra 필요성 실험 유료 실행 시점 질문에 답) \intent{설계·구현 후 확인받고 (Recommended)}
  \item \intent{지금 10만원 정도는 비용 한도 책정이 가능해}
  \item \chg{(09-25 정본 §82) Astra는 Astra만 할 수 있는 일에만: J1 과제 컴파일(계약 + 물체 지식·검출 이름 후보·위험·감사 지점·새로움 표시·예상 실패별 대응, 분포 안 과제는 계약 캐시), J2 실패 진단 + 복구안(판정은 V1h·코드), J3 검사 술어 합성, J4 새 물체 이름·속성(위치 찾기는 SAM 3.1), J5 저빈도 진행 감사(하트비트 대체), J6 오프라인 경험 증류. 강등: 5 s 하트비트 ack·단계 경계 ack·성공 판정·정체 첫 대응 $\rightarrow$ 작은 모델·코드(V1h \tentative{0.949·0.44 ms} 대 Astra 첫 토큰 약 3 s). 호출 = 결정이 바뀔 수 있을 때만(T0/T\_fail/T\_nov/T\_audit), 성공 편 \tentative{1--3회}(지금 K2 \tentative{7--12회}). effort low·high 모두 논문에. 유료 한도 약 10만 원, 실행은 구현 뒤 승인. 구현은 R7 관문 뒤 $\rightarrow$ 지금 런타임은 K2(5 s 하트비트) = 기준선}
  \item \dirn{(§82 E-Astra-necessity) 같은 위층 자리에 Astra low/high, Qwen3-VL-32B·8B·4B, 위층 없음, 무작위 시점 Astra. 층: 분포 안(대조)·새 물체·긴 과제·새 실패·random. 오프라인 슬롯 $\rightarrow$ 폐루프. 판정 초안 $\Delta$ = Astra $-$ 최고 로컬 $\ge$ \tentative{+10 pt}·하한 $>0$ $\rightarrow$ 필수, 비열등이면 강등. 논문 계획 절 표(\texttt{tab:astra}), 결과 칸 주황}
  \item \intent{나는 그냥 아스트라 자체가 이 움직임에도 개입을 했으면 좋겠어 근데 딜레이가 문제니까 이걸 해결하면서 도 개입을 하는 방법을 찾는거야 최소화 하는 것도 방법이겠지 VLA가 한거를 감사하는건 아스트라의 능력중 극히 일부일 뿐인거구}
  \item (제안 A: M4 느린 투표자 / B: 물체 기준 경로 / C: 조건부 동작 프로그램에 답) \intent{A,B 더하는 방식이 조금 더 좋은 것 같긴하다}
  \item \intent{근데 아스트라가 경로를 짜는 방식으로 물체를 잡아 단독으로 할때? 아니면 직접 보면서 엔드포인트 보는 방식으로 조종을 해서 그다음에 움직이지 않나?}
  \item \intent{아니 근데 아스트라가 만약에 3D 공간상에 예측해서 경로를 그려주는 일을 잘한다면 그 방법으로 시도를 해보는 것도 좋을 수는 있을 것 같아}
  \item \intent{탐침 E-Astra-motion: 3D 경로 그리기 대 보고 조종하기 이거 한번 해보자용}
  \item \intent{뎁스를가지고 한다고 하면 될 것 같아 뎁스가 있잖아 ai워커가}
  \item \intent{아니다 뎁스를 쓰면 너무 범용성이 떨어질 것 같은 생각이드네 이미지 상만으로 어떻게 3D를 방향을 제대로 계획을 시킬 수 있을까...?}
  \item \intent{그냥 뎁스라기보다는 그 가우시안 스플레팅인가? 아니면 포인트 클라우딩 포인트 클라우딩으이었던거 같은데 스테레오캠으로 하는 그거 그거를 어떻게 잘 해서 안되는건가? 공간 좌표 만들고 거기서 하는거.?}
  \item \dirn{(user-log 78·79 브레인스토밍, \tentative{탐침 준비 중, 결과 없음}) 움직임 인터페이스 A+B: (A) Astra의 typed 답을 M4 원장에 늦게 온 표로 넣어 합의·실행 뒤 확인을 거쳐야만 확정, (B) Astra가 물체 기준 경로를 그리고 조이스틱 VLA가 따라가며 순간 보정. E-Astra-motion 탐침 = 3D 경로 그리기 대 보고 조종하기. 깊이는 범용성이 떨어져 RGB만으로 3D: 지지 평면 광선 교차, 두 시점 삼각측량, 스테레오 점군 + 가상 시점. 논문 방법 ``움직임 인터페이스''·계획 절에 주황으로}
  \item \intent{아게 진짜로 잘 모르긴하겠다 아스트라를 어케 쓰지?}
  \item \idea{(Astra 감독자 틀 정리) 과제 전 = J1(흐름·계약·예상 실패 대응) / 과제 중 = 움직임 개입(A+B, \tentative{시험 중})·J5 저빈도 감사 / 실패 때 = J2 진단·복구안 / 과제 뒤 = J6 오프라인 교사(작은 모델 규칙·재학습 라벨). 5 s 하트비트는 현재 구현 K2 = E-M8c 기준선, K5(감사 + 모를 때만)를 새 사전 등록으로 추가}
\end{itemize}

\subsection{M9 실패 복구}
\begin{itemize}
  \item \intent{VLA에서 가장 좋은 실패 복구 방법을 찾는다. 논문은 무조건 1년 반 이내.}
  \item \intent{VLA에서 실패했을 때 어떻게 되돌아갈지 정하는 프레임워크를 참고한다.}
  \item \chg{v3의 ``Astra는 맨 마지막'' 폐기 $\rightarrow$ 실패 판정 순간 Astra를 비동기로 바로 부르고, 기다리는 동안 아래 층(스킬 재시도 $\rightarrow$ Jev 선택)이 로봇을 움직임}
  \item \chg{FLARE~\cite{flare2026}는 자세 오류/환경 파괴 구분만 가져옴(재시도는 학습형이었음)}
  \item \dirn{복구 = 언제 / 어디로 / 무엇을 남길지. 어디로 = 사전조건이 다 참인 가장 늦은 체크포인트. Jev가 코드가 낸 복구 제안 \{아래 층 결과 인정, 체크포인트 재개, 리셋, 안전 대기, 모르겠음(올림)\} 중 고름(FAIL은 뒤집지 못함). Astra는 실패 순간 이미 비동기로 불려 있음 (로봇 FaRe, LLM RIR이 같은 달 같은 분해)}
  \item \idea{복구 기법끼리 직접 비교한 벤치마크는 아직 없음 $\rightarrow$ ``최고''는 수치로 못 정함}
  \item \idea{(09-24) Harness VLA는 동작마다 결과를 진행/회복 가능/불가 셋으로 분류(우리 M7 판정과 대응). 단 플래너가 결과 파일을 기다리는 동기 방식}
  \item \chg{반대 증거 추가: CaP-X에서 검증·디버깅을 강하게 시킨 프롬프트가 오히려 약간 나빴음(68.29$\rightarrow$65.43) $\rightarrow$ ``반성 유지''는 잠정 그대로, 비교 조건 필요}
\end{itemize}

\subsection{M10 경험 축적}
\begin{itemize}
  \item \intent{경험 축적이 제대로 되어야 한다.}
  \item \intent{API로 불러오는 모델의 경험 축적을 어떻게 하는지, 어떻게 해야 효율적인지 논문을 찾는다. 1년 반 이내.}
  \item \idea{Astra: 단순 경험 검색 + 교훈~\cite{ouyang2025reasoningbank,zhang2025ace}. Jev: 교훈 목록 말고 딱 맞는 규칙 0--2개만(MemCompiler~\cite{memcompiler2026}에서 얻은 착상, 심사 전)}
  \item \chg{``메모리 없는 게 낫다''는 반대 논문이 실재~\cite{budget2026memory} $\rightarrow$ 같은 예산 비교 필수}
  \item \dirn{(09-24) 교훈 항목에 적용 조건·증상·증거 수·신뢰도·``반증 조건'' 칸 추가(RPent 메모리 형식). 여러 셀에서 확인된 것만 Astra에 넣는 안}
  \item \idea{Harness VLA는 메모리가 없으면 위치 바꾼 과제에서 87$\rightarrow$31로 떨어짐 $\rightarrow$ 과제별 기록이 성능의 큰 몫. 단 Astra low 92.63\%는 논문이 아니라 저장소 리더보드 값(LIBERO-PRO T+S 8칸, seed 1--10, 메모리 스냅샷 157/200 + 584/600, 에피소드당 약 412 s 멈춰서 생각)}
\end{itemize}

\subsection{단계 2 설계 1판 --- 모듈마다 한 줄 (2026-09-23)}
\begin{itemize}
  \item \idea{공통: 실패 판정은 M7 한 곳, 멈춤은 첫 계획과 복구 중 안전 대기뿐, 확률은 보정 재기 전엔 안 씀}
  \item M1: \dirn{술어 사전 하나(코드 기하 / 임계 민감 / VLM 필요 3등급). 변환 방법 기본은 후보 A(09-24, Claude가 정함)}
  \item M2: \dirn{Astra가 ``세션 계약''을 씀: 목표·가정·단계·결정 지점(스킬에 박힌 id 중에서 고름)·금지 조건}
  \item M3: \dirn{보기는 서로 뚜렷한 몇 개만(보기가 많고 비슷하면 정확도 떨어짐). 사용자 원안(촘촘) 대 혼합안 비교}
  \item M4: \dirn{겹친 답을 ``표''로 모음. 동시통역의 앞부분 합의 + 추측 디코딩의 유예 창 + 실제 움직임 확인}
  \item M5: \dirn{같은 보기끼리만 RTC식으로 섞고 마지막에 저크 제한(Ruckig, 잠정 기본)}
  \item M6: \dirn{스킬 카드(LLM 도구 스키마식) + 단계별 Jev 질문 틀. ``기존 스킬'' = 생성형 / 보유 라이브러리 둘 다 \decide{}}
  \item M7: \dirn{하드(코드 기하 술어만) + 소프트 채널 두 개 이상이면 실패}
  \item M8: \dirn{``이때쯤 나왔어야'' 두 구현 비교. 모델 스스로 판단은 늦게 울릴 것이라는 가설(BAGEN 예산 자기 추정 낙관, 트리거 시점을 잰 것 아님). JEV-Star식 주기+사건 호출·REFLEX~\cite{reflex2026}식 확신 기반 올려 보내기는 비교 조건 후보(아직 안 넣음). 격자 10프레임 기본 유지}
  \item M9: \dirn{실패 순간 Astra 비동기 호출, 그동안 아래 층이 움직임. 되돌아갈 곳 = 조건이 다 맞는 가장 늦은 체크포인트}
  \item M10: \dirn{Astra엔 사례+교훈, Jev엔 딱 맞는 규칙 0--2개만(없으면 안 넣음). 다 쌓으면 오히려 나빠짐}
  \item \idea{로봇 밖에서 가져온 것: 복구 = 데이터베이스 사가(되돌릴 수 있는 단계만 보상 동작으로 되돌림), 스무딩 = 게임 애니메이션 관성화, 스킬 = 타입 있는 빈칸 검사, 계약 = 형식 검증}
\end{itemize}

\subsection{융합 VLA와 학습 단계 (09-24 -- 09-25)}
\begin{itemize}
  \item \intent{어떤식으로 vlm을 써야할지를 좀 생각해봐야할듯 엔드투 엔드 목적으로 사용하게 될거니까 하긋ㅂ할때 어케하게될지}
  \item \chg{(09-24 §51·§52, 선택 답 ``A$\rightarrow$B 단계적'') C 영점 $\rightarrow$ A 결정층만 학습(보기 NLL SFT $\rightarrow$ 스냅샷 분기 기대 보상 + KL RL) $\rightarrow$ B 같은 백본에 action expert(flow matching, KI 기울기 차단). 보정 = 질문별 온도 + conformal 문턱, 모델 바뀔 때마다}
  \item \chg{(09-24 §53·§54) 영점 선택기는 최빈 기준선보다 낮음(0.48 대 0.67) $\rightarrow$ 폐루프엔 학습한 결정층만. 옛 정답은 관측으로 정의되지 않음(코드 규칙 상한 0.80--0.92) $\rightarrow$ labels\_v2(현재 손가락 중점 $\rightarrow$ 소단계 목표까지 남은 변위, 1 mm 상태에서 관문 통과)}
  \item \intent{그 손목캠의 정보도 들어가는건 별로 일까? 좌우 손목캠이 항상들어가는것도 좋은데 실제 상호작용하는게 활성화되어서 보게하는}
  \item \dirn{(09-25 §57) 머리캠은 항상 + 지금 스킬을 실행 중인 팔의 손목캠(D405 424$\times$240). 양팔 스킬일 때만 두 손목. 활성 팔은 계약의 \texttt{arms}로 정해짐. 네 조건 절제로 확정}
  \item \intent{어떻게 여러 이미지를 하나에 처리가 되게 할지를 알아봐야겠듯}
  \item \dirn{(09-25 §59, D27) 타일로 합치지 않고 Qwen3-VL 기본 다중 이미지(머리 252 + 손목 104 토큰, 표지 문구). 병목은 토큰 수가 아니라 같은 이미지 요청의 동시 진입 $\rightarrow$ lead 호출(첫 시퀀스 먼저) + 접두·멀티모달 캐시 필수. 설계 p95 0.28 s}
  \item \intent{각 모듈별로 따로 놀면서ㅗ 가각의 성능에 의존하는 파이프라인보다는 최대한 하나로 융합되는게 좋은 연구인거알지?}
  \item \intent{어쩔수없는건 물론 그냥 하긴해야하구}
  \item \chg{(09-25 §58) 런타임 = VLA 하나(Qwen3-VL-4B): 이미지 $\rightarrow$ typed 결정 토큰(트라이 재정규화 확률, M4가 소비) + action expert 연속 행동 + 보조 기하 헤드(학습 신호). 모델 밖 = Astra, M4, 안전 투영. 명시적 인식·스크립트 스킬은 모듈형 기준선이자 교사}
  \item \decide{주 표의 주 행을 융합 모델로 둘지(현재 융합·모듈형 두 행 병기). 융합 런타임에서 M4 (b)·critic의 물체 술어를 무엇으로 확인할지}
  \item \intent{일단 지금은 나란히 하고 나중에 확실히 정리가 되면 융합으로 해야지근데 그거와 별계로 저 M4의 크리틱 저거는 무조건 해결을해 1년반 내 논문 싹다 깃도 뒤져서 신뢰성있는걸로 찾아서 해봐봐}
  \item \chg{(09-24 §60·§61·§64, 위 [결정 필요] 둘 다 답) 주 표는 지금 융합·모듈형 병기, 결과로 정리되면 융합으로. M4 (b)·critic 측정원은 둘로 나눔: 로봇 쪽 = 고유 감각 코드 T1(그리퍼 열림·쥠/빈 집기·쥔 채 들림, 하드 채널), 세계 쪽 = 같은 백본의 확인 헤드 V1h(\texttt{on}·\texttt{in\_contact}·\texttt{near}·접촉 멈춤, 소프트 채널, 술어별 온도 + conformal, 빈 집합 = \texttt{unknown}). critic 세계 쪽 경보 = V1h 단독. 근거(1.5년 안, 신뢰도 높음): Self-Improving EFM(같은 모델 진행 예측이 성공 판정기로 더 안정), TA-VLA(모터 전류 토크가 말단 상태 구분), SAFE(토큰 불확실성 기준선 약함)}
  \item \idea{(E-M4b-meas, 파이프라인 확인이라 수치는 임시) V1h \tentative{0.949} 대 명시 인식 \tentative{0.596}·영점 \tentative{0.750}, T1 \tentative{0.992--0.998}, V1h + 토크 충돌 규칙 재현율 \tentative{0.786} $<$ V1h 단독 \tentative{0.871} $\rightarrow$ 토크 규칙은 명목 운동 토크에 걸려 문턱을 망침. 명시 인식은 3D 오차 \tentative{35\,mm}라 측정원 탈락}
  \item \chg{(09-24 §67 C4) ``모델 한 번 호출'' $\rightarrow$ 스텝마다 두 호출: decide(공유 접두 순전파 1회 $\rightarrow$ 결정 확률 + 문맥 은닉 + 확인 헤드) $\rightarrow$ M4 $\rightarrow$ chunk(캐시 문맥 + 확정 결정 $\rightarrow$ expert, CUDA 그래프). 백본 순전파는 스텝당 1회라 융합 원칙 그대로. 융합 런타임은 HF 백본 서버, vLLM은 모듈형 기준선 전용}
  \item \chg{(09-24 §65) 결과 기반 라벨: 사전 등록 규칙이 \texttt{plan}을 골랐지만 판별력 없음 $\rightarrow$ 학습 정답 아님, 실패 보기 거부권 용도만. 학습 정답은 labels\_v2 유지}
  \item \dirn{(09-25 user-log 71, 계획·\tentative{검증 중}) 결정 입력에 시간 맥락: 머리·활성 손목 카메라마다 $[t-0.3\,\text{s},\,t]$ 2프레임을 Qwen3-VL 비디오로(시각 토큰 수 불변). 동기 = S-E2E 진단(정확도 \tentative{약 0.72}, 오답은 움직임 유무·크기 한 칸, 부호 반대 \tentative{8/155}, 한 시점 입력엔 속도가 없음). chunk는 decide 캐시 문맥을 그대로 읽음. 채택은 E-TC 2$\times$2 결과로(실험 절 참고)}
  \item \chg{(09-25 13:05 UTC E-TC 사전 등록, 학습 전) 입력 정의 확정: $\Delta=0.3$\,s = 10\,Hz 프레임 $k-3$(에피소드 시작은 0으로 자름). 실제 처리기 탐침에서 정지 이미지와 2프레임 클립 모두 머리 \tentative{252}·손목 \tentative{104} 토큰(grid\_t 1), 우리 [과거, 현재] 채움 픽셀 = 공식 비디오 처리기 출력 비트 동일. 공식 비디오 경로의 타임스탬프 \tentative{+8} 텍스트 토큰 대신 \texttt{<|image\_pad|>} 유지 + 카메라 표지 ``head camera (2 frames: t-0.3 s, t):''(프롬프트 해시 파일 무수정). 움직임 줄 = \texttt{motion: arm=<still|slow|fast> gripper=<closing|still|opening>}, 속도는 인과 후방 차분(행의 \texttt{qd}는 중앙 차분이라 0.1\,s 미래를 읽음 $\rightarrow$ 안 씀), 팔 3분위 \tentative{0.218·0.607}\,rad/s, 그리퍼 0.85 분위 \tentative{0.176}/s, 줄 드롭아웃 p \tentative{0.3}(값만 \texttt{unknown}). 과거 프레임 드롭아웃·$\Delta$ 0.1\,s는 이번 등록 밖. 판본 표지 \texttt{D27v2-video2}·\texttt{se2e-motion@v1}}
  \item \chg{(09-25 정본 §83, user-log 80) E-TC 사전 등록 판정대로 결정 입력에 움직임 줄(\texttt{se2e-motion@v1}) 채택, 2프레임 비디오 불채택(CPU 전처리로 p95 증가). 개요·모델 그림의 주황 점선 과거 프레임을 지우고 텍스트에 ``움직임 줄''. 런타임 적용(새 직렬화 판본, 런타임 속도 구간화, 보정·카나리 재생성, R2 줄 생성) = \tentative{다음 작업}}
  \item \idea{(§83 발견) S-E2E 변환 행의 \texttt{proprio.qd}·그리퍼 속도 = 중앙 차분(0.1 s 미래 포함). 움직임 줄은 안 쓰지만 action expert 고유감각 입력엔 들어감 $\rightarrow$ 변환기를 인과 후방 차분으로 고친 새 데이터 판본, 확인 실험 \tentative{진행 중}. 지금까지 S-E2E 판정은 §56 파이프라인 시험으로 유지 + 한계 명시}
\end{itemize}

\subsection{M4 사전 등록 적합화 (R7 순회, 09-25)}
\begin{itemize}
  \item \dirn{원칙: 사전 등록은 구속한다. 코드는 사전 등록 문장을 그대로, 사전 등록이 열어 둔 값은 정본에 근거 한 줄과 함께}
  \item \chg{(§72) 유예 창 $W$ = 첫 도전 표 뒤에 더 필요한 같은 도전 표 수. $W=0$은 첫 도전 표에서 즉시 교체, 기본 $W=1$(도전 표 2개), 비가역 $W{+}1$}
  \item \chg{(§73) $\tau=1$(순서형 이웃 1칸), 접촉 근처(목표 $\le$ 5\,cm 또는 접촉)는 $\tau=0$. 경계 직후 $W$ 2·$\gamma$ 1.0, 큐 임계 조기 호출 $g$, 장면 변화 조기 호출, 미확정 실행 (b) 엄격화는 사전 등록 조건에 없음 $\rightarrow$ SCOPED(절제 후보)}
  \item \chg{(§74) $\gamma$ = 정확히 2/3(``3표 중 2'', 정수 곱 비교), 런타임과 E0.5 재생이 같은 함수. C2 = Slow Brain VLM Stream 원문 그대로(동시 요청 상한 없음, 매 틱 요청 시각이 가장 새 유효 응답, 5\,s 타임아웃 $\rightarrow$ 기본 행동). Holm은 ``하한 $>0$'' 방향만(판정 2·카나리), 판정 10은 양측. E-M4-lat의 STALE\_MAX = 1.5\,s $\times d^\ast_{p95}/d_{p95}$(E0)(비례 규칙, 지연 큐 구현 때)}
  \item \chg{(§75·§76) $H=1$ = 같은 스텝을 시각을 당겨 여러 번 묻기: 송신 $t$의 호출은 시작 시각이 $[t+\hat d,\,t+\text{lead}_{\max}]$인 모든 스텝의 표. $\text{lead}_{\max}$ \tentative{1.0\,s}(E0 뒤 런타임 스텝당 표 중앙값 $\ge 2$ 규칙으로 재결정, 유한 양수만). 판정 7($H=3$ 대 $H=1$)은 다스텝 DecCall 전까지 SCOPED. 나이 비교: 정확히 1.5\,s·5\,s는 유지}
  \item \dirn{(R7 12회차 D3) (b) 범주를 다음 결정 요청 상태 텍스트에 ``\texttt{last\_step: <범주>}'' 한 줄로 되먹임 --- 사전 등록 C5의 정의. 없으면 런타임 C5 = C5$'$라 판정 3(C5 대 C5$'$)이 성립하지 않음. 구현됨(§77, 아래 줄; 결정 호출 형식이 바뀌어 학습 형식도 함께 바꿈)}
  \item \idea{(R7 12회차) 그 밖의 결함: E0.5 FLIP\_TH 후보 분위가 $\alpha=0.01$이 아님, E0.5 지표 2개 미보고, 호출 기록에 질문별 확률 분포·요청/응답 원문·Astra 원응답 없음 $\rightarrow$ 수정됨(§77, 아래 줄)}
  \item \chg{(§77 D3) 모든 결정 요청 상태의 마지막 줄 = \texttt{last\_step: <none|OK|LAG|DEVIATE|CONTRADICT>}. C4·C5·C6 = 마지막으로 끝난 스텝의 (b) 범주(늦은 V1h DEVIATE는 OK·LAG를 DEVIATE로 올림), C0--C3·C5$'$ = \texttt{none}. (§79 N1 정정) 결정 호출은 경계 판정 뒤에 만들어 경계 틱의 호출부터 방금 끝난 스텝의 범주가 실림(DEVIATE·CONTRADICT의 조기 호출은 다음 틱). 행동 expert 문맥엔 안 넣음. 직렬화 판본 ser-A-min-2 $\rightarrow$ 옛 체크포인트·보정 파일 거부(전부 §56 파이프라인 시험이라 논문 수치 손실 없음). 학습 자료 줄 = 그 편이 기록한 실행 뒤 확인(R2는 이미 기록된 직전 스텝 확인이라 생성기 변경 없음)}
  \item \chg{(§77 D4) 호출 기록: 질문별 보기 확률(보정 전후), \texttt{last\_step}, 요청 본문·백엔드 원응답은 내용 해시 이름 파일(\texttt{<run>/blobs/<sha256>})에 한 번만. Astra = 원응답 텍스트·최대 토큰·effort·머리캠 JPEG 원본. 결정 호출 이미지는 해시만(원본이면 약 18\,MB/편)}
  \item \chg{(§77 D1·D2·N6·N7) FLIP\_TH 후보 q0.99(사전 등록 $\alpha$ 0.01)·q0.999·q0.95, E0.5 같은 시각 뒤집힘을 성공 궤적·섭동 궤적으로 나눠 보고 + 이름 절제 판 A1--A4의 A0 대비 뒤집힘, 카나리 표류 의심이면 그 전 날짜 보정 파일로는 확률 게이트(J5) 끔, E0 판정 4(스텝당 표 중앙값 $\ge 2$) 집계 구현}
  \item \idea{(§77 무작위 대조) 가짜 세계 160판: 행동·원장·호출 수는 구현 전과 같고 요청 마지막 줄만 다름 $\rightarrow$ 차이는 줄을 읽는 학습 모델에서만 생김(판정 3이 재는 것)}
  \item \chg{(§77 보충, 메인 결정) C5$'$는 줄을 지우지 않고 \texttt{none}(지우면 판정 3이 (b) 정보가 아니라 입력 형식 변화를 잼). (b) 범주 희소성(R2 DEV 결정 스냅샷 \tentative{1,081개 중 DEVIATE 21, CONTRADICT 0})은 본 단계 B 학습 사전 등록 때 비 OK 가중·섭동 설계로 정함. 성능이 부족하면 user-log 70대로 1.5년·신뢰도 규칙의 선행 기법을 찾아 새 판·새 사전 등록으로 적용}
  \item \idea{(R7 13회차 D1--D3) 12회차 수정이 인용한 사전 등록 문장을 전체로 다시 대조하니: FLIP\_TH 초기값은 질문 평균 점수의 선형 보간 분위였음 $\rightarrow$ 사전 등록은 \texttt{question\_id@vN}별 split conformal(표본이 모자라면 문턱 없음). 카나리 기준일은 질문 판본을 안 봐 ser-A-min-2 판이 옛 기준과 비교됨 $\rightarrow$ 판본 바뀌면 새 기준일. 표류 의심 다음 날 깨끗한 카나리만으로 J5가 다시 켜짐 $\rightarrow$ 표류 의심 뒤엔 그날 이후 다시 맞춘 보정 파일이 있어야 켬. 수정됨(§79, 아래 줄)}
  \item \chg{(§79, 09-25 11:50 커밋 2c4a30d) FLIP\_TH = 질문 판본별 split conformal: 성공 스텝 점수의 $\lceil (n+1)(1-\alpha)\rceil$번째 값, 순위 $>n$이면 문턱 없음(None, $\alpha$ 0.01 $\rightarrow$ 성공 스텝 99개 필요; 13회차 모의 판 n 47 $\rightarrow$ 없음). 카나리 기준일 = 같은 세트 $\times$ 같은 \texttt{question\_id@vN} 집합의 첫 카나리(없으면 그날이 새 기준일, 옛 판본과 비교 안 함). J5 = 그 모델 지문의 가장 최근 표류 의심 날짜 $\le$ 보정 파일 날짜일 때만 켬(다음 날 깨끗한 카나리만으론 안 켬). 경계 틱 호출의 \texttt{last\_step} = 방금 끝난 스텝(가짜 세계 160판 중 \tentative{156판} 비트 동일, 달라진 4판은 범주가 바뀐 경계 틱에 주기 호출이 걸린 판, 성공 여부 모두 같음). 빈 입력 거부, Isaac 워커 \texttt{OMP\_WAIT\_POLICY=PASSIVE}}
  \item \dirn{(§79 보충, 메인 결정) 런타임 \texttt{flip\_th}는 지금 질문 공통 한 값(끔) $\rightarrow$ J4를 켜기 전(E0.5 뒤 FLIP\_TH 확정 때) 질문별 사전으로 바꾸는 것이 필수 작업. E0.5 점수 창(직전 3표)과 런타임 창(최근 4표 대 그 앞 4표)도 그때 한쪽으로 맞춤}
  \item \idea{(R7 14회차, 2c4a30d 대상) \tentative{FAIL}: 결함 1 --- 결과 기반 라벨을 읽는 코드(평가 \texttt{--truth outcome:}, 단계 A 결과 라벨)가 §78 (1) 라벨 신뢰 기준(재생 비트 동일)을 적용하지 않아 무효 886행이 정답에 섞일 수 있음. 문서 2(옛 §77 (ii) 문장, 새 빌드 29편 재라벨 계획). 13회차 수정은 모두 동작으로 확인, 하드 리셋 뒤 같은 워커 C5 $\rightarrow$ C5$'$ 행동 \tentative{1,592개} 비트 동일. 연속 무결 \tentative{0}, 수정 \tentative{진행 중}}
  \item \chg{(§80, 커밋 0790dd4) 14회차 수정 완료: 결과 라벨을 \emph{읽는} 모든 곳(평가 \texttt{--truth outcome:} $\rightarrow$ \texttt{e05}·\texttt{rd}·\texttt{calib}, 단계 A 결과 라벨, 자체 점검 대조)이 한 함수 \texttt{replay\_bit\_identical}(재생 편차가 수이고 정확히 0)로 거름. 편차 기록 없음·null도 불신(기준 = 확인된 비트 동일). 단위는 행. 뺀 수를 결과에 기록(풀 \tentative{886 / 6,627행} 제외). S-E2E 판정 스크립트 입력 검사(끝 3점 = 사전 등록 평가 일정, 재개 창 2001--2050, 구동기 출력 필수 $\rightarrow$ 판정 통과 불변), 빈 \texttt{--seeds} 선택 거부}
  \item \idea{(R7 15회차, 9cee78b 대상) \tentative{FAIL}: 결함 0, 문서 2 --- (1) 고정 선행 실행이 재실행 복원을 비트 동일하게 만든다는 옛 채택 서술(\texttt{pool.md} 등)에 §78 정정 표시 없음, (2) 자체 점검 불신 597행은 ``13편''이 아니라 13시드·20편. 14회차 수정은 실제 풀 라벨로 동작 확인(불신 886행·29편이 정답에 0), 새 시드 DEV 12에서 하드 리셋 뒤 같은 워커 C5 $\rightarrow$ C5$'$ 행동 \tentative{1,877개} 비트 동일. 연속 무결 \tentative{0}, 문서 정정 \tentative{진행 중}}
\end{itemize}

\subsection{Astra--VLA 결합: 직렬 흐름과 일반 조종 (09-25 16:17--18:02 UTC, user-log 83--88)}
\begin{itemize}
  \item \intent{로우로하자 그럼 이 로우와 하위 vla가 rtc로서 서로서로 보완하면 좋을듯 아스트라는 반응성이 느려서 그러지 vla는 방향성이 빠르지만 자기가 뭘하는지에대한 성능이 조금 떨어져서 그렇지}
  \item \intent{vla가 못하지는 않은데, 아스트라보다 덜하지 뭘하는지는 대충알지만 실시간성이 좋다는거지}
  \item \chg{(정본 §82 보충 2) 실행 중 Astra effort = low, high는 흐름에 안 씀(첫 토큰 45--73\,s). Astra = 느리지만 무엇을 하는지 앎, VLA = 대략 알지만 실시간 $\rightarrow$ RTC식 상호 보완}
  \item \intent{이게 뭐랄까 엄청 촘촘하게 이루어질 수 있어야하는데 api를 직렬이 아니라 1초단위로 3개 불러오면 1초1초가 연속적인 방법이 될 수는 있잖어}
  \item \intent{API의 문제를 vla가 vla의 문제를 api아스트라가 해결해줘야해 api는 이걸 흐름으로 아는게 아니니까 매번 다른답을 주는데(물론 그걸 최소화하도록 해야겠지)}
  \item \intent{목표를 공간에 표현하는 것 조차 안하고 그냥 일반적 아스트라 조종법을 따르겠다는 얘기}
  \item \intent{네, 일반 조종법만}
  \item \chg{(설계 개정 16:33) 공간 목표(3D 점·경로·이미지 점) 없음. Astra = GPT-as-Policy식 일반 조종: 진행 평가(progress, progressing/failed/uncertain, aligned/misaligned, 확신, 근거 시점) 뒤 continue / edit(말단 $\le$5\,cm·$\le$0.35\,rad·그리퍼) / stop. 흐름 고정: 직전 답 대비 차이 + 근거 필수(F1) 대 매번 새로(F0) $\rightarrow$ 탐침}
  \item \intent{아스트라의 조정을 부드럽게 해결하는 방식이어야한다는 뜻인거지 뭐 그리고 반응성이 너무 느리고 이게 이미지만 보고는 확실하지 않은 경우들이 많기 때문에 그렇지}
  \item \chg{(설계 §11) edit = 말단 기준 편향, 1--3\,s 선형 램프·감쇠, 한 답 50\% / 연속 두 답 같은 방향 100\% / 반전이면 0. 속도 0.08\,m/s·가속 0.32\,m/s$^2$ 상한. uncertain·확신 low $\rightarrow$ continue(불확실만으론 안 가로챔), ``잠시 느리게 더 보기''만 허용. VLA가 1\,s 넘게 반대면 편향 절반 + 다음 요청에 표시}
  \item \intent{아스트라는 3개의 캠을 다봐야해 예를들어 헤드캠에서 봤을때는 원시때문에 물체를 잡을 때, 잡았다고 오탐하지 않게 우론목에는 안보이는 그런게 있는거지 그걸 해결할 수 이썽야하는거지}
  \item \chg{(설계 §12) Astra = 머리 + 두 손목(이름표). 접촉·잡음·놓음 주장은 그 팔 손목 시점 근거 필수(코드가 확인), 어긋나면 손목 우선 + T1 교차. 비가역 동작은 두 층 + 손목/T1 근거}
  \item \idea{(탐침 무료 사전 실행, 중간 관찰) Qwen3-VL-8B는 세 대를 볼 때 오히려 ``잡았다'' 과잉 주장 \tentative{12/14}, 머리만이면 실제 파지 \tentative{5/5}를 놓침 $\rightarrow$ 카메라를 늘리는 게 늘 좋지 않음, 규칙은 코드로 강제하고 구성은 G1로 정함}
  \item \chg{(정본 §86, 탐침 결과로 확정) 요청 형식 \textbf{F0}(F1 뒤집힘 감소 \tentative{0.594$\rightarrow$0.381, 36\%} $<$ 기준 50\%, 답 쌍 적음). 늦은 답의 edit 폐기 문턱 6\,s $\rightarrow$ \tentative{15\,s}(흐름 요청 p50 \tentative{8.9--10.0\,s}, 답 나이 p95 \tentative{11.1--14.2\,s}; 6\,s면 \tentative{70--100\%} 버림). \textbf{파지 확정 = T1·V1h}, Astra 영상 판정은 참고(손목 포함해도 실제 쥠 \tentative{13/25} 놓침, 테 건 쥠 \tentative{6/6}, high 표적 교정 \tentative{2/8}; 머리만이면 uncertain \tentative{39/40}). 확신 low 답 \tentative{0회} $\rightarrow$ 끼어듦은 합의·속도 상한·VLA 확인으로 거름. 부드러운 반영 저크 \tentative{0.83} 대 즉시 \tentative{2.36\,m/s$^3$}}
  \item \idea{(과잉 주장의 원인) Qwen 세 대 \tentative{12/14}는 카메라 수가 아니라 프롬프트에 ``쥠'' 정의가 없던 탓(정의 뒤 \tentative{1/7}, 왼손목 가려도 \tentative{6/7}) $\rightarrow$ 모든 판정 질문 프롬프트는 대상의 조작적 정의를 담는다}
  \item \intent{VLA도 사진을 보는것도 좋을 것 같구}
  \item \chg{E-CAM3(예정): VLA 입력 머리 + 활성 손목 대 머리 + 두 손목(시각 토큰 356$\rightarrow$460), E-TC 채택 틀}
  \item \intent{엔드포인트의 움직임을 위에 덧그리는 방식도 뭔가 도움이 될 수 있을 것 같기도 하구}
  \item \chg{(설계 §13) Astra 영상에 현재 손끝, 최근 2--3\,s 궤적(흐려짐), 다음 청크 예정 화살표, 직전 편향, 로봇 축 표시. 예정 화살표는 VLA 입력엔 안 넣음(정답 누수). 궤적은 MolmoAct 0--255 형식으로 기록}
  \item \intent{몰모 엑트를 많이 참조분해해서 잘사용 적용해보면 좋을듯}
  \item \idea{(MolmoAct 분해, 연구 문서 molmoact\_deepdive) 깊이 토큰(RGB에서 단안 추정 $\rightarrow$ VQ-VAE 100코드) $\rightarrow$ 이미지 궤적(1--5점) $\rightarrow$ 행동, 매 스텝 \tentative{200--260}토큰 생성, 사슬 단계 절제 없음. MolmoAct2는 궤적 생성 버림, 깊이 경로 \tentative{0.787\,s} 대 기본 \tentative{0.179\,s}}
  \item \dirn{가져올 것(실행 비용 0만): (1) 미래 손끝 궤적 보조 손실, (2) 그린 명령을 따르도록 학습(되짚기 명령 절반), (3) expert 층별 KV. 안 가져올 것: 실행 중 깊이·궤적 생성, 가중치, 데이터 혼합}
  \item \intent{1}\textcolor{red}{\textasciitilde}\intent{3 다 허용할게 몰모엑트 많이 참고하고}
  \item \intent{1,2,3 허용했지만 옳은 방향으로 바꿔도돼 확인해보고}
  \item \chg{(정본 §84·보충 1) 결합 설계 승인, E-MA1 승인, falsify 먼저 평가(같은 서명 재사용 한 번, 부분 일치 재사용 안 함, J5·J6에 남김). 재점검: 1\,Hz 조밀 호출은 비용 한도로 불가 $\rightarrow$ 접촉 구간 위주, 답 나이 상한 \tentative{6\,s}, E-MA1 해석 주의(궤적 보조 = 라벨과 같은 정보)}
  \item \intent{그냥 아스트라 부르는거 그냥 1번씩만 하는걸로하자}
  \item \chg{(정본 §84 보충 2) 흐름 호출 = 동시 요청 1개, 답이 오면 최신 관측으로 다음 요청(간격 = 지연 $L$). 계단식 겹침 폐기. 합의까지 약 $2L$ --- 그사이는 VLA가 메움}
  \item \chg{(구현 계획 16과제, 메인 판정) 두 층 합의의 뜻 = 신선한 Astra 답(3\,s 이내)이 반대하지 않고 T1 전제가 참이면 비가역 동작 진행(긍정 답을 기다리면 파지마다 수 초 정지). stop은 기록·critic 신호만, 로봇을 세우지 않음}
\end{itemize}
